%% ══════════════════════════════════════════════════════════════
%%  CHAPTER 4 — FULL PIPELINE RECONSTRUCTION
%% ══════════════════════════════════════════════════════════════
\chapter{Paper Generation Pipeline}
\label{ch:pipeline}

\section{Pipeline Overview}
\label{sec:pipeline:overview}

The paper generation pipeline is the central execution path of \autoscholar, transforming a research prompt into a published, peer-reviewed academic paper.
The entire process executes within a single streaming HTTP connection lasting 45--90 minutes, involving up to 180 recursive tool calls, and consuming 500K--1M tokens of Claude context.

This chapter reconstructs the pipeline phase by phase, documenting the code involved, data flows, decision points, and failure modes at each stage.

\section{Phase 1: Trigger and Initialization}
\label{sec:pipeline:trigger}

\subsection{Trigger Mechanisms}

The pipeline can be triggered through three mechanisms:

\begin{description}
  \item[Scheduled generation.] The \file{autoPaperScheduler.ts} (1,285 lines) runs on a configurable cadence (default: 1 hour). The scheduler maintains a cursor into the 1,000+ entry prompt bank and selects the next prompt.

  \item[Admin-triggered generation.] Administrators can trigger generation via \texttt{POST /api/admin/trigger-paper} (random prompt) or \texttt{POST /api/admin/trigger-paper-custom} (specified prompt).

  \item[Interactive generation.] Users submit research prompts through the web interface, which may be detected as paper-generation requests and upgraded.
\end{description}

\subsection{Prompt Selection}

For scheduled generation, the prompt selection algorithm works as follows:

\begin{algorithm}[H]
\caption{Prompt Selection in the Auto-Paper Scheduler}
\label{alg:prompt-selection}
\SetAlgoLined
\KwIn{Prompt bank $\mathcal{P}$ with 1,000+ entries, cursor position $c$, selection mode $m$}
\KwOut{Selected prompt entry $p$}

\uIf{$m = \text{sequential}$}{
  $p \gets \mathcal{P}[c \bmod |\mathcal{P}|]$\;
  $c \gets c + 1$\;
}
\uElseIf{$m = \text{weighted\_random}$}{
  $w_i \gets 1.0$ for all $i$\;
  \ForEach{recently used domain $d$}{
    $w_i \gets w_i \times 0.5$ for all $p_i$ with domain $d$\;
  }
  $p \gets \text{WeightedSample}(\mathcal{P}, w)$\;
}
\Else{
  $p \gets \text{UniformRandom}(\mathcal{P})$\;
}

\Return{$p$ with fields: domain, method, topic, full prompt text}\;
\end{algorithm}

The prompt bank is constructed deterministically at import time from three sources:
\begin{itemize}
  \item \textbf{24 domains} --- 18 finance (market microstructure, volatility modeling, asset allocation, ...) + 6 stats/ML/AI domains.
  \item \textbf{24 methods} --- 18 econometric (panel GMM, event studies, DCC-GARCH, ...) + 6 stats/ML/AI methods.
  \item \textbf{14 topics per domain} --- Specific, publishable research questions.
  \item \textbf{Compatibility matrix} --- Maps each domain to its compatible methods, preventing nonsensical pairings.
\end{itemize}

The iteration order (variation, topic, domain) ensures consecutive prompts always switch domains, providing topic diversity.

\subsection{Surprise Prompt Generation}

For interactive sessions, the ``surprise paper'' feature (\func{generateSurprisePrompt}) picks a random prompt from the full pool, weighted to avoid recently selected ones (tracked in localStorage with a 30-entry cap).

\subsection{Pipeline Run Initialization}

When a scheduled run begins:

\begin{enumerate}
  \item A \texttt{PaperRun} record is created in the database with status \texttt{queued}.
  \item A session ID is generated for the conversation.
  \item The \governor is initialized with default state (confidence 1.0, all retry budgets full).
  \item A heartbeat timer starts (10-second interval) to detect stalled runs.
  \item The run status transitions to \texttt{researching}.
\end{enumerate}

\section{Phase 2: Research Prompt Enhancement}
\label{sec:pipeline:enhancement}

\subsection{Research Detection}

Before reaching Claude, the prompt passes through \func{isResearchPaperRequest} (\file{researchPromptEnhancer.ts}), which uses two-part regex matching:

\begin{itemize}
  \item \textbf{Verb patterns:} ``write'', ``generate'', ``create'', ``produce'', ``draft'', ``compose''
  \item \textbf{Noun patterns:} ``paper'', ``article'', ``study'', ``thesis'', ``report'', ``analysis''
\end{itemize}

\subsection{Methodology Detection}

\func{detectResearchMethodology} classifies the prompt into one of four tracks:

\begin{itemize}
  \item \textbf{Econometrics:} Keywords include ``regression'', ``panel'', ``GMM'', ``GARCH'', ``causal'', ``instrumental''.
  \item \textbf{Statistics:} Keywords include ``hypothesis'', ``bootstrap'', ``Bayesian'', ``distribution'', ``non-parametric''.
  \item \textbf{Machine Learning:} Keywords include ``random forest'', ``XGBoost'', ``neural'', ``classification'', ``prediction''.
  \item \textbf{AI:} Keywords include ``deep learning'', ``transformer'', ``NLP'', ``attention'', ``fine-tuning''.
\end{itemize}

\subsection{Prompt Upgrading}

\func{upgradeResearchPrompt} calls Claude Opus non-streaming (temperature 1.0 for creative variety) to transform vague prompts into detailed 400--700 word research specifications.
Each methodology type receives tailored guidance:

\begin{description}
  \item[Econometrics:] Must include identification strategies, 6--10 specific variable definitions, and endogeneity considerations.
  \item[Statistics:] Must specify distributional assumptions, test statistics, and power analysis considerations.
  \item[Machine Learning:] Must name a baseline + 4 models with cross-validation strategy and hyperparameter search space.
  \item[AI:] Must specify architecture details, training regime, and ablation study design.
\end{description}

The upgrader preserves the language of the original prompt (supports French and English).

\section{Phase 3: Literature Review (Gauss Phase)}
\label{sec:pipeline:literature}

\subsection{System Prompt Activation}

When the upgraded prompt reaches Claude, the 1,798-line system prompt activates the \gauss role for the first phase.
The system prompt instructs:

\begin{quote}
\itshape
``Begin with Agent Gauss. Search at least 3 academic databases. Identify 15--30 relevant papers. Extract methods, datasets, and limitations. Identify 3--5 research gaps.''
\end{quote}

\subsection{Execution Sequence}

\gauss typically executes the following tool sequence:

\begin{enumerate}
  \item \texttt{academic\_search\_multi} --- Parallel search across Semantic Scholar, OpenAlex, arXiv, CrossRef, and Google Scholar.
  \item \texttt{literature\_gap\_finder} --- If the initial search returns sufficient results, run the full 7-step gap analysis pipeline.
  \item Individual database queries for targeted follow-up.
  \item \texttt{web\_search\_exa} or \texttt{tavily\_search} --- For recent working papers not yet indexed in academic databases.
\end{enumerate}

\subsection{Outputs}

The \gauss phase produces:
\begin{itemize}
  \item A list of 15--30 relevant papers with BibTeX entries.
  \item Identified research gaps ranked by novelty, feasibility, and impact.
  \item A refined research question based on gap analysis.
  \item Suggested methodology and data requirements.
\end{itemize}

\section{Phase 4: Data Acquisition (Turing Phase)}
\label{sec:pipeline:data}

\subsection{Data Strategy}

Based on \gauss's findings, \euler transitions to the \turing role and plans a data acquisition strategy.
The system prompt defines strict data source hierarchies with fallback chains:

\begin{enumerate}
  \item For stock/company data: FMP $\to$ EODHD $\to$ web scraping.
  \item For macro data: FRED $\to$ World Bank $\to$ IMF $\to$ OECD.
  \item For academic datasets: structured APIs $\to$ \gls{swdes} web extraction.
  \item For alternative data: Polymarket $\to$ Google Trends $\to$ news sentiment.
\end{enumerate}

\subsection{Typical Tool Sequences}

A typical \turing phase for a cross-sectional asset pricing study might involve:

\begin{enumerate}
  \item \texttt{get\_historical\_price} for multiple assets (portfolio returns).
  \item \texttt{get\_income\_statement} / \texttt{get\_balance\_sheet} for fundamental factors.
  \item \texttt{get\_treasury\_rates} for the risk-free rate.
  \item \texttt{download\_fmp\_data} to export bulk data to CSV.
  \item \texttt{get\_wb\_indicator\_data} for country-level controls.
  \item \texttt{execute\_custom\_python\_analysis} to merge and validate datasets.
\end{enumerate}

\subsection{Data Provenance}

The system prompt mandates data provenance documentation:
\begin{quote}
\itshape
``For every dataset used, record: source API, endpoint, parameters, date range, number of observations, and any transformations applied.''
\end{quote}

This information is later embedded in the paper's data section and the project ZIP's metadata files.

\section{Phase 5: Data Engineering (Newton Phase)}
\label{sec:pipeline:engineering}

\subsection{Execution}

\newton primarily operates through \texttt{execute\_custom\_python\_analysis}, running code that:

\begin{enumerate}
  \item Loads downloaded CSV/JSON files into pandas DataFrames.
  \item Cleans and standardizes column names, date formats, and numeric types.
  \item Merges datasets by date or entity identifiers.
  \item Constructs derived variables (returns, ratios, dummies, interactions).
  \item Runs preliminary diagnostics (summary statistics, correlation matrices, stationarity tests).
  \item Saves the prepared dataset for the analysis phase.
\end{enumerate}

\subsection{Script Management}

Each Python execution is tracked by the ScriptManager (\file{scriptManager.ts}, 276 lines):

\begin{itemize}
  \item Scripts are saved to session-scoped temp directories with unique IDs.
  \item Up to 8 repair attempts are allowed per script.
  \item The diffEngine supports unified diffs, search-replace blocks, and fuzzy matching.
  \item Structured errors are parsed and formatted for Claude's next turn.
\end{itemize}

\section{Phase 6: Empirical Analysis (Fisher Phase)}
\label{sec:pipeline:analysis}

\subsection{Model Selection Logic}

\fisher's model selection follows the methodology track detected during prompt enhancement:

\begin{description}
  \item[Econometrics track:] Start with OLS baseline, add fixed effects, instrument with 2SLS if endogeneity is suspected, robustness checks with alternative specifications, clustering of standard errors.
  \item[Statistics track:] Formulate null/alternative hypotheses, select appropriate test statistics, compute p-values and confidence intervals, assess power.
  \item[ML track:] Train baseline (logistic regression or OLS), train 4 advanced models, cross-validate, compare via appropriate metrics (RMSE, AUC, F1), compute SHAP values for interpretation.
  \item[AI track:] Implement architecture, train with early stopping, run ablation studies, visualize attention patterns.
\end{description}

\subsection{Robustness Checks}

The system prompt mandates robustness checks:

\begin{enumerate}
  \item Alternative variable definitions (e.g., different return windows).
  \item Subsample analysis (pre/post crisis, by sector, by size).
  \item Alternative model specifications (e.g., Newey-West vs. clustered SE).
  \item Placebo tests where applicable.
  \item Sensitivity to outlier removal.
\end{enumerate}

\subsection{Figure Generation}

Each analysis produces figures following the academic figure style directive.
The \governor tracks figure count and enforces limits based on the active profile (8--15 figures depending on error state).
Generated figures pass through:

\begin{enumerate}
  \item Matplotlib rendering with JF-style configuration.
  \item Blank figure detection (axes artist inspection + pixel variance analysis).
  \item AutoScholar watermark addition.
  \item Base64 encoding and SSE transmission to the client.
  \item Optional vision analysis via Claude Opus for figure description.
\end{enumerate}

\section{Phase 7: Paper Writing (Euler Write Phase)}
\label{sec:pipeline:writing}

\subsection{Section-by-Section Generation}

\euler writes the paper one section at a time using the \texttt{write\_paper\_section} tool.
The tool definition specifies extremely detailed length and quality requirements, targeting 60--90 page papers.

\subsubsection{Empirical Paper Sections}

\begin{enumerate}
  \item \textbf{Introduction} (untitled, following JF convention) --- Motivation, contribution, preview of results.
  \item \textbf{Literature Review} --- Organized by theme, positions the paper relative to existing work.
  \item \textbf{Methodology} --- Formal model specification, estimation strategy, identification.
  \item \textbf{Data} --- Sources, construction, summary statistics, preliminary analysis.
  \item \textbf{Results} --- Main findings, robustness checks, subsample analysis.
  \item \textbf{Discussion} --- Economic interpretation, implications, comparison with existing literature.
  \item \textbf{Conclusion} --- Summary, contributions, limitations, future research.
\end{enumerate}

\subsubsection{Theoretical Paper Sections}

\begin{enumerate}
  \item \textbf{Economic Environment} --- Model setup, agents, preferences, technology.
  \item \textbf{Equilibrium} --- Equilibrium concept, characterization, existence/uniqueness.
  \item \textbf{Computational Method} --- Numerical solution, calibration, simulation.
  \item \textbf{Quantitative Results} --- Model implications, moment matching, counterfactuals.
  \item \textbf{Extensions} --- Model variants, sensitivity analysis.
\end{enumerate}

\subsection{Writing Style Enforcement}

During the writing phase, the \governor injects the JF Writing Style Directive (282 lines), which includes:

\begin{itemize}
  \item 100+ banned words/phrases organized by severity tier.
  \item Section-specific conventions (e.g., Introduction is unnumbered, abstract $\leq$ 100 words).
  \item Citation formatting rules (``Author (year)'' for narrative, ``(Author, year)'' for parenthetical).
  \item Advanced prose techniques: emphatic stress positioning, topic-comment sentence structure.
  \item AI detection countermeasures: lexical diversity targets, syntactic variation requirements.
  \item An 8-pass verification checklist for self-review.
\end{itemize}

\subsection{Section Rejection Tracking}

The tool executor (\file{toolExecutor.ts}) tracks section rejection to prevent infinite loops.
If a section is rejected more than 3 times (due to insufficient length or quality), the system accepts the best available version and moves on.

\section{Phase 8: LaTeX Compilation}
\label{sec:pipeline:compilation}

\subsection{Section Assembly}

The \texttt{compile\_research\_paper} tool (\file{latex/compiler.ts}, 91 lines) sorts sections into canonical academic order using a predefined ordering map:

\begin{lstlisting}[style=typescript, caption=Section ordering map.]
const SECTION_ORDER = {
  'introduction': 1,
  'literature_review': 2,
  'methodology': 3,
  'data': 4,
  'results': 5,
  'discussion': 6,
  'conclusion': 7,
  'appendix': 100,
};
\end{lstlisting}

Appendix sections automatically receive a \texttt{\textbackslash appendix} command.

\subsection{Template Selection}

Three document templates are available (\file{latex/paperGenerator.ts}, 1,235 lines):

\begin{description}
  \item[\texttt{article}] --- Standard two-column academic article with \texttt{natbib} bibliographic style.
  \item[\texttt{working\_paper}] --- Single-column working paper format with wider margins.
  \item[\texttt{econometrica}] --- Mimics the \emph{Econometrica} journal style.
\end{description}

All templates share a common preamble with 50+ \LaTeX{} packages and custom macros.

\subsection{Compilation Pipeline}

\begin{algorithm}[H]
\caption{\LaTeX{} Compilation Pipeline}
\label{alg:latex-compilation}
\SetAlgoLined
\KwIn{Sections, BibTeX, figures, template}
\KwOut{Compiled PDF path}

Write \texttt{paper.tex}, \texttt{references.bib}, and figure files to temp directory\;
Sanitize \LaTeX{} content (escape special chars, fix Greek letters, handle math mode)\;
Auto-fit wide tables (wrap in \texttt{adjustbox} or \texttt{resizebox})\;

\textbf{Pass 1:} Run \texttt{pdflatex -interaction=nonstopmode paper.tex}\;
\If{BibTeX present}{
  \textbf{Pass 2:} Run \texttt{bibtex paper}\;
  \textbf{Pass 3:} Run \texttt{pdflatex} (resolve references)\;
  \textbf{Pass 4:} Run \texttt{pdflatex} (finalize)\;
}

\If{compilation failed}{
  Register compilation with \texttt{latexPatchManager}\;
  \For{attempt $\gets 1$ \KwTo $8$}{
    Parse errors via \texttt{latexLogParser}\;
    Apply diff patch via \texttt{diffEngine}\;
    Recompile\;
    \If{success}{\textbf{break}\;}
  }
}

\Return{PDF path or error}\;
\end{algorithm}

\subsection{LaTeX Sanitization}

The sanitizer (\file{latex/sanitizer.ts}, 670 lines) fixes common Claude-generated \LaTeX{} issues:

\begin{itemize}
  \item Unescaped special characters (\%, \&, \$, \#, \_).
  \item Greek letter handling outside math mode.
  \item Mismatched math delimiters.
  \item Em-dash and en-dash normalization.
  \item Wide table detection and automatic wrapping.
  \item Landscape environment stripping (known to cause compilation issues).
  \item Blank PNG figure detection (decompresses IDAT chunks and measures pixel variance).
\end{itemize}

\subsection{Iterative Repair}

The \LaTeX{} patch manager (\file{latex/latexPatchManager.ts}, 369 lines) provides a read-patch-recompile cycle:

\begin{enumerate}
  \item \texttt{inspectLatexErrors} parses the log file, classifying errors into 8 types (undefined control sequence, math, unicode, environment, package, fatal, missing character, other).
  \item \texttt{patchLatex} applies diffs using the shared diffEngine with fuzzy matching.
  \item \texttt{recompileLatex} runs the full pdflatex+bibtex pipeline.
  \item Up to 8 repair attempts are allowed before giving up.
\end{enumerate}

\section{Phase 9: Packaging}
\label{sec:pipeline:packaging}

The \file{latex/packaging.ts} (358 lines) creates a ZIP archive following the \gls{rre} directory structure:

\begin{verbatim}
project/
  paper.tex              % Main LaTeX source
  references.bib         % BibTeX references
  outputs/
    paper.pdf            % Compiled PDF
    figures/             % All generated figures
  code/
    01_data_download.py  % Pipeline-ordered scripts
    02_data_prep.py
    03_analysis.py
    ...
  data/
    raw/api/             % API-sourced data
    raw/web/             % Web-extracted data
    processed/           % Cleaned datasets
    metadata/            % Data provenance files
  run_pipeline.py        % Auto-generated master script
  requirements.txt       % Auto-detected dependencies
  README.md              % Build instructions and documentation
\end{verbatim}

Key features:
\begin{itemize}
  \item Scripts are ordered by pipeline stage using numeric prefixes.
  \item Data files are routed to structured subdirectories based on source type.
  \item \texttt{requirements.txt} is auto-generated by parsing imports from all scripts and mapping to pip package names.
  \item \texttt{run\_pipeline.py} is auto-generated with sequential script execution.
  \item \texttt{README.md} includes abstract, build instructions, project structure, and data provenance.
\end{itemize}

\section{Phase 10: Validation (Popper Phase)}
\label{sec:pipeline:validation}

\subsection{Review Trigger}

After successful compilation, the scheduler invokes \popper:

\begin{enumerate}
  \item The run status transitions to \texttt{validating}.
  \item The full paper content (title, abstract, sections, BibTeX) is submitted to \func{generatePeerReview}.
  \item Claude reviews the paper at temperature 0.3 for consistency.
  \item The review result includes:
  \begin{itemize}
    \item Overall score (0--100).
    \item Recommendation (accept/minor\_revision/major\_revision/reject).
    \item 6-dimension detailed scores.
    \item Section-by-section feedback.
    \item Strengths and weaknesses lists.
  \end{itemize}
\end{enumerate}

\subsection{Publication Decision}

The \governor evaluates the Popper review via \func{evaluatePopper}:

\begin{itemize}
  \item Score $\geq 45$: Paper is published (run status $\to$ \texttt{published}).
  \item Score $< 45$: Paper is rejected (run status $\to$ \texttt{rejected}).
  \item The review is stored in the \texttt{peer\_reviews} table.
\end{itemize}

\section{Phase 11: Publication}
\label{sec:pipeline:publication}

\subsection{ASID Assignment}

Published papers receive an \gls{asid} via \file{doiService.ts} (195 lines):

\begin{equation}
\text{ASID} = \texttt{ASID-}\langle\text{year}\rangle\texttt{-}\langle\text{sequence}\rangle
\end{equation}

where sequence is a zero-padded 4-digit number generated by atomic database counter increments (thread-safe for concurrent generation).

\subsection{Storage}

The published paper record includes:
\begin{itemize}
  \item Title, abstract, domain classification, template.
  \item Full \LaTeX{} sections (stored as JSON).
  \item BibTeX references.
  \item Figure count, word count, section count, \LaTeX{} line count.
  \item Peer review score and recommendation.
  \item Project ZIP file path.
  \item ASID.
  \item Associated session ID for cost tracking.
\end{itemize}

\subsection{Domain Classification}

The domain detector (\file{utils/domainDetector.ts}) classifies papers into domains based on keyword analysis of the title and abstract, using a weighted scoring system across 24 domain categories.

\section{Checkpoint and Resume System}
\label{sec:pipeline:checkpoint}

\subsection{Motivation}

Given that pipeline runs last 45--90 minutes, crash recovery is essential.
The checkpoint system allows runs to resume from the last saved state after server restarts.

\subsection{Checkpoint Structure}

\begin{lstlisting}[style=typescript, caption=Pipeline checkpoint structure.]
interface PipelineCheckpoint {
  runId: string;
  sessionId: string;
  prompt: string;
  messages: ClaudeMessage[];
  governorState: AppState;
  decisionLog: DecisionLogEntry[];
  toolResults: ToolResult[];
  writtenSections: { name: string; content: string }[];
  figureCount: number;
  stage: PipelineStage;
}
\end{lstlisting}

\subsection{Checkpoint Lifecycle}

\begin{enumerate}
  \item Checkpoints are saved every 8 seconds during active generation.
  \item On server restart, the scheduler queries for runs with status \texttt{researching}, \texttt{analyzing}, \texttt{drafting}, or \texttt{validating}.
  \item Active runs are marked as \texttt{failed} (to prevent zombie runs).
  \item Administrators can manually resume from a checkpoint via the admin API.
  \item Resumed runs restore the full conversation history, governor state, and continue from the last tool call.
\end{enumerate}

\section{Pipeline Metrics}
\label{sec:pipeline:metrics}

Typical pipeline execution metrics:

\begin{table}[H]
\centering
\caption{Typical pipeline execution metrics.}
\label{tab:pipeline-metrics}
\begin{tabular}{lr}
\toprule
\textbf{Metric} & \textbf{Typical Value} \\
\midrule
Total duration & 45--90 minutes \\
Tool calls & 80--180 \\
Python executions & 15--40 \\
Figures generated & 8--15 \\
Context tokens consumed & 500K--1M \\
API cost per paper & \$15--\$40 (Opus 4.6 pricing) \\
Paper length & 60--90 pages \\
BibTeX references & 30--80 \\
\LaTeX{} compilation attempts & 1--4 \\
Popper review score & 50--85 (mean $\approx$ 65) \\
\bottomrule
\end{tabular}
\end{table}
